\documentclass[aspectratio=169]{beamer}

\usetheme{Madrid}
\definecolor{unsblue}{RGB}{30,60,120}
\definecolor{unslightblue}{RGB}{220,230,245}
\definecolor{alertred}{RGB}{180,30,30}
\usecolortheme[named=unsblue]{structure}
\setbeamercolor{palette primary}{bg=unsblue, fg=white}
\setbeamercolor{palette secondary}{bg=unsblue!80, fg=white}
\setbeamercolor{palette tertiary}{bg=unsblue!60, fg=white}
\setbeamercolor{palette quaternary}{bg=unsblue, fg=white}
\setbeamercolor{block title}{bg=unsblue, fg=white}
\setbeamercolor{block body}{bg=unslightblue}
\setbeamercolor{block title alerted}{bg=alertred, fg=white}
\setbeamercolor{block body alerted}{bg=red!8}
\setbeamercolor{block title example}{bg=unsblue!60!black, fg=white}
\setbeamercolor{block body example}{bg=unsblue!8}

\usepackage[utf8]{inputenc}
\usepackage[T1]{fontenc}
\usepackage{babel}
\babelprovide[main, import]{spanish}
\usepackage{amsmath, amssymb}
\usepackage{bm}
\usepackage{booktabs}

\newcommand{\R}{\mathbb{R}}
\newcommand{\bx}{\bm{x}}

\title[MAD1 -- Unidad 4]{Aprendizaje Supervisado}
\subtitle{Clase 21: Ensambles y Random Forest}
\author[J. Bavio]{Dr.\ José Bavio}
\institute[UNS]{Departamento de Matemática\\Universidad Nacional del Sur}
\date{2026}

\begin{document}

\begin{frame}
  \titlepage
\end{frame}

\begin{frame}{Contenido de la clase}
  \tableofcontents
\end{frame}

% ═════════════════════════════════════════════════════════════════════════════
\section{La idea de los ensambles}
% ═════════════════════════════════════════════════════════════════════════════

\begin{frame}{¿Por qué combinar modelos?}
  \begin{block}{Problema}
    Un árbol de decisión tiene alta varianza: cambiás un poco los datos y el árbol cambia drásticamente.
  \end{block}

  \medskip
  \begin{exampleblock}{Intuición: la sabiduría de la multitud}
    Si le preguntás a una sola persona cuántos caramelos hay en un frasco, puede errar por mucho.\\
    Si promediás las respuestas de 100 personas independientes, el promedio suele ser muy preciso.\\[4pt]
    Los \textbf{métodos de ensamble} hacen exactamente eso con modelos.
  \end{exampleblock}

  \medskip
  Condición clave: los modelos deben ser \textbf{suficientemente distintos entre sí}.\\
  Promediar 100 copias del mismo modelo no ayuda nada.
\end{frame}

% ─────────────────────────────────────────────────────────────────────────────
\section{Bagging}
% ─────────────────────────────────────────────────────────────────────────────

\begin{frame}{Bootstrap}
  Antes de entender bagging, necesitamos el \textbf{bootstrap}.

  \medskip
  \begin{block}{Muestra bootstrap}
    Dado un conjunto $\mathcal{D}$ de $n$ observaciones, una muestra bootstrap
    $\mathcal{D}^{(b)}$ es una muestra de tamaño $n$ extraída \textbf{con reemplazo} de $\mathcal{D}$.
  \end{block}

  \medskip
  \begin{exampleblock}{Coloquialmente}
    Es como sacar bolitas de una bolsa y volver a meterlas antes del próximo sorteo.
    Algunas bolitas van a salir dos veces, otras ninguna.
    En promedio, una muestra bootstrap contiene el $\approx 63.2\%$ de las observaciones originales (distintas).
  \end{exampleblock}

  \medskip
  Las observaciones que no entran en la muestra bootstrap ($\approx 36.8\%$)
  se llaman \textbf{out-of-bag} (OOB). Ya veremos por qué eso es útil.
\end{frame}

\begin{frame}{Bagging (\textit{Bootstrap Aggregating})}
  \textbf{Algoritmo} (Breiman, 1996):

  \medskip
  \begin{enumerate}
    \item Para $b = 1, \ldots, B$:
    \begin{enumerate}
      \item Generar una muestra bootstrap $\mathcal{D}^{(b)}$.
      \item Entrenar un árbol $T^{(b)}$ sobre $\mathcal{D}^{(b)}$.
    \end{enumerate}
    \item Combinar por \textbf{mayoría de votos}:
    \[
      \hat{y}(\bx) = \arg\max_k \sum_{b=1}^{B} \mathbf{1}\bigl(\hat{y}^{(b)}(\bx) = k\bigr)
    \]
  \end{enumerate}

  \medskip
  \begin{exampleblock}{¿Por qué reduce la varianza?}
    Si un modelo tiene varianza $\sigma^2$ y los modelos son independientes,
    el promedio de $B$ tiene varianza $\sigma^2/B$.
    La varianza cae con $B$. El sesgo queda igual.
  \end{exampleblock}
\end{frame}

% ─────────────────────────────────────────────────────────────────────────────
\section{Random Forest}
% ─────────────────────────────────────────────────────────────────────────────

\begin{frame}{El problema del bagging puro}
  \begin{alertblock}{Problema: correlación entre árboles}
    Si hay un feature muy predictivo (por ejemplo, una variable que domina la predicción),
    \textbf{todos} los árboles van a usarla en los primeros splits.\\
    Resultado: los $B$ árboles son muy parecidos entre sí.\\
    Promediarlos no reduce mucho la varianza.
  \end{alertblock}

  \medskip
  \begin{exampleblock}{Analogía}
    Preguntarle a 100 personas que leyeron exactamente el mismo artículo no es
    más informativo que preguntarle a una sola.
    Para que el promedio ayude, las fuentes deben ser \textbf{distintas}.
  \end{exampleblock}
\end{frame}

\begin{frame}{Random Forest: decorrelación de árboles}
  \begin{block}{Ingrediente extra de Random Forest (Breiman, 2001)}
    En cada split de cada árbol, se elige el mejor split \textbf{solo entre un subconjunto aleatorio} de $m$ features (no los $p$ totales).
  \end{block}

  \medskip
  \begin{itemize}
    \item Clasificación: $m = \lfloor\sqrt{p}\rfloor$ (por defecto)
    \item Regresión: $m = \lfloor p/3 \rfloor$ (por defecto)
  \end{itemize}

  \medskip
  \begin{exampleblock}{¿Por qué funciona?}
    Al forzar que cada árbol ignore la mayoría de los features en cada split,
    los árboles se ``ven obligados'' a ser distintos entre sí.
    Cuando se promedian árboles poco correlacionados, la varianza cae mucho más.
  \end{exampleblock}
\end{frame}

\begin{frame}{Varianza del promedio: por qué la decorrelación importa}
  Sean $B$ árboles, cada uno con varianza $\sigma^2$ y correlación $\rho$ entre pares.
  La varianza del promedio es:
  \[
    \rho\,\sigma^2 + \frac{1-\rho}{B}\,\sigma^2
  \]

  \medskip
  \begin{itemize}
    \item Si $B \to \infty$: la varianza tiende a $\rho\,\sigma^2$.
    \item Si $\rho = 0$ (árboles independientes): varianza $\to 0$.
    \item Si $\rho = 1$ (árboles idénticos): varianza $= \sigma^2$ (sin mejora).
  \end{itemize}

  \medskip
  \begin{exampleblock}{Conclusión práctica}
    Aumentar $B$ (más árboles) siempre ayuda, pero con rendimientos decrecientes.
    Reducir $\rho$ (decorrelación vía aleatoriedad en features) es lo que hace único a RF.
  \end{exampleblock}
\end{frame}

% ─────────────────────────────────────────────────────────────────────────────
\section{Estimación Out-of-Bag}
% ─────────────────────────────────────────────────────────────────────────────

\begin{frame}{Error Out-of-Bag (OOB)}
  \begin{block}{Idea}
    Cada árbol $T^{(b)}$ se entrenó sobre $\mathcal{D}^{(b)}$ que excluye
    $\approx 36.8\%$ de los datos. Esas observaciones son las \textbf{OOB} de ese árbol.
    Para cada $\bx_i$, predecir usando \textbf{solo} los árboles que no vieron $\bx_i$.
  \end{block}

  \medskip
  \[
    \hat{y}_i^{\text{OOB}} = \arg\max_k \sum_{b:\, i \in \text{OOB}(b)} \mathbf{1}(\hat{y}^{(b)}(\bx_i) = k)
  \]

  \medskip
  \begin{exampleblock}{¿Para qué sirve?}
    El error OOB es un estimador \textbf{casi insesgado del error de generalización}
    y no requiere reservar un conjunto de validación separado.\\[4pt]
    Es como tener validación cruzada ``gratis'' mientras entrenás.
  \end{exampleblock}
\end{frame}

% ─────────────────────────────────────────────────────────────────────────────
\section{Importancia de Variables}
% ─────────────────────────────────────────────────────────────────────────────

\begin{frame}{Importancia por impureza (MDI)}
  \begin{block}{Mean Decrease in Impurity (MDI)}
    La importancia de la variable $j$ es la reducción total de impureza (Gini/entropía)
    acumulada en todos los splits de todos los árboles que usan $j$, promediada sobre $B$.
  \end{block}

  \medskip
  \begin{exampleblock}{Coloquialmente}
    Una variable es ``importante'' según MDI si, cada vez que se la usa para hacer un split,
    los nodos hijos quedan mucho más puros que el padre.\\
    Si una variable hace splits que no mejoran casi nada la pureza, su importancia es baja.
  \end{exampleblock}

  \medskip
  \begin{alertblock}{Cuidado}
    MDI \textbf{sobreestima} la importancia de variables con alta cardinalidad
    (muchos valores posibles, como IDs o fechas).
    Es un sesgo conocido del método.
  \end{alertblock}
\end{frame}

\begin{frame}{Importancia por permutación (MDA)}
  \begin{block}{Mean Decrease in Accuracy (MDA)}
    Para cada árbol y cada variable $j$:
    \begin{enumerate}
      \item Calcular el error OOB base.
      \item \textbf{Permutar aleatoriamente} los valores de $x_j$ en el conjunto OOB.
      \item Recalcular el error con los valores permutados.
      \item Importancia de $j$ = aumento promedio del error.
    \end{enumerate}
  \end{block}

  \medskip
  \begin{exampleblock}{Intuición}
    Si mezclar al azar los valores de una variable arruina las predicciones,
    esa variable era importante.\\
    Si el modelo no se entera de que la mezclamos, era ruido.\\[4pt]
    Es como preguntarse: ``¿qué pasa si le saco esta pieza al motor?''
  \end{exampleblock}
\end{frame}

% ─────────────────────────────────────────────────────────────────────────────
\section{Hiperparámetros y Práctica}
% ─────────────────────────────────────────────────────────────────────────────

\begin{frame}{Hiperparámetros principales de Random Forest}
  \begin{center}
  \begin{tabular}{lll}
    \toprule
    \textbf{Hiperparámetro} & \textbf{Efecto} & \textbf{Valor típico} \\
    \midrule
    \texttt{n\_estimators}   & Número de árboles            & 100--500 \\
    \texttt{max\_features}   & Features por split           & \texttt{'sqrt'} \\
    \texttt{max\_depth}      & Prof.\ máxima del árbol      & None \\
    \texttt{min\_samples\_leaf} & Mínimo muestras en hoja  & 1--5 \\
    \texttt{bootstrap}       & Usar muestras bootstrap      & True \\
    \texttt{oob\_score}      & Calcular error OOB           & True \\
    \bottomrule
  \end{tabular}
  \end{center}

  \medskip
  \begin{block}{Regla práctica}
    Más árboles ($n\_\text{estimators}$) casi siempre mejora o mantiene igual el rendimiento:
    no existe ``sobreajuste por tener demasiados árboles''.
    El límite es computacional, no de precisión.
  \end{block}
\end{frame}

\begin{frame}{Ventajas y limitaciones de Random Forest}
  \begin{columns}
    \column{0.5\textwidth}
    \begin{block}{Ventajas}
      \begin{itemize}
        \item Muy robusto: difícil hacerlo sobreajustar
        \item Buena estimación de error OOB
        \item Importancia de variables incorporada
        \item Pocas decisiones de diseño
        \item Paralelizable
      \end{itemize}
    \end{block}

    \column{0.5\textwidth}
    \begin{alertblock}{Limitaciones}
      \begin{itemize}
        \item Caja negra: no podés explicar una predicción puntual fácilmente
        \item Lento para predecir en tiempo real (muchos árboles)
        \item Memoria: guarda $B$ árboles completos
        \item No extrapola bien fuera del rango de entrenamiento
      \end{itemize}
    \end{alertblock}
  \end{columns}
\end{frame}

\begin{frame}{Resumen de la clase}
  \begin{itemize}
    \item Los \textbf{métodos de ensamble} combinan múltiples modelos débiles para obtener uno robusto.
    \item \textbf{Bagging} reduce la varianza entrenando sobre muestras bootstrap y promediando.
    \item \textbf{Random Forest} agrega aleatoriedad en la selección de features para \textbf{decorrelacionar} los árboles.
    \item El \textbf{error OOB} es una estimación gratuita del error de generalización.
    \item La \textbf{importancia de variables} (MDI y MDA) permite interpretar el modelo de forma global.
  \end{itemize}

  \bigskip
  \textbf{Próxima clase:} SVM y KNN --- dos enfoques muy distintos al problema de clasificación.
\end{frame}

\end{document}
